\scanpage{54}

Algèbres étales et la différente (3)\hfill Tribu n° 65

sur $A$ tel que, en posant $B'=A[x]$, $F'(x)$ soit inversible dans $B'_{q\cap B'}$ et $B'_{q\cap B'}\to B_q$ soit bijectif. Généralisation au cas où $B$ est localisée \unsure{et une|d'une} $A$-algèbre finie.

Si, de plus, $A_{q\cap A}$ est normal et si $B$ est étale en $q$, $B_q$ est normal. La prop. 4 (p. 4) du texte.

Ajouter, entre astérisques, les améliorations dépendant du main-theorem.

\uline{n° 5.}\quad \uline{Groupes d'automorphismes.}

La prop. 6 du texte, dans le cas noethérien non réduit (cf. Tribu n° 64). Son cor. (p. 7) est la prop. appliquée à $H$, de sorte qu'on le vide.

\uline{n° 6.}\quad \uline{Un critère de non ramification.}

Il s'agit de l'actuel n° 4 (§ 1) amélioré comme suit. Remplacer (ici et partout) l'idiote formule $(A^*)^{-1}B=L$ par $KB=L$. Référer à Alg. IX, § 2 pour la définition du discriminant.

Dans le b) du th. 1, ajouter

(iii)\quad $B$ est finie sur $A$ et $B/m.B$ est étale sur $A/m$ ;

(iv)\quad $B$ est étale finie sur $A$ ;

comme conditions équivalentes à (i) et (ii). Ça absorbe la partie intéressante de c), ainsi que le cor. 1.

Pour démontrer le lemme (p. 10--11) on remarque que $P(x)=D(1,x,\ldots,x^{n-1})$ est fonction polynôme de $x$, et qu'on a $P\neq0$ (passer à la clôture algébrique). On a donc $x\in R$ tel que $P(x)\neq0$.

Le cor. 1 est absorbé. La déf. 2 est supprimée. Le cor. 2 devient une prop. donnant l'équivalence de :

a)\quad $B$ est étale finie sur $A$ ;

b)\quad $B_m$ est étale finie sur $A_m$ pour tout idéal maximal $m$ de $A$ ;

c)\quad $B$ est finie sur $A$ et $B/m.B$ est étale sur $A/m$ pour tout idéal maximal $m$ de $A$.

Le cor. 3 est vidé. On donne enfin le cor. 4 (y remplacer ``non-ramifié'' par ``étale'').

On essaiera de mettre l'\uline{appendice} suivant : si $I$ est un idéal nilpotent de $A$, il y a une équivalence de catégories entre les algèbres étales finies sur $A$ et les algèbres étales finies sur $A/I$. Voir S. G. A.

\mspage{53}
